\chapter{Tests and Conclusion}

\section{Tests}

\subsection{Test Environment}

Both the agent have been tested along the building process continously,
to find bugs, test different configurations, and improve them.
Most of the test have been performed with the server and pddl running locally, this has been
done following the instructions on the relative repositories:\dots

\subsection{Test Scenarios}

The decision was to test multiple times on the scenarios of Challenge 1 and Challenge 2.
In particular for the second part we checked that the instruciton messages sent by the Admin
were satisfied by the agent and registered and try to fix eventual problems.

\subsection{Results and Limitations}

The most difficult part has been the prompring of the LLM agent,
to have it doing the tasks correctly and usign the right tools.

One known limitaiton is that it is not always able to parse multiple revards for multiple requests
in one message, and therefore if the first reward is negative for example, it doesn't execute the instructions.
Another problem that has been noted is that on multiple questions, more than
3 in one message, if there is an evaluate\_math\_expression call the model puts its output at the end.
It was difficult to spot the possible pattern that the model extracts from the prompt
in order to convince itself that the correct way to output is this one. So this 2 problems are still open
and mostly related to the prompt.

\section{Conclusion}

Summarize the main contributions and discuss future developments.
